\chapter{לקוחות מוקדמים והתאמת מוצר-שוק}

התאמת מוצר-שוק — \term{Product--Market Fit (PMF)} — היא הרגע שבו המוצר פותר בעיה
אמיתית עבור שוק מוגדר, באופן שהשוק ``מושך'' את המוצר מידי הצוות. מארק אנדריסן
תיאר זאת כתחושה שאי אפשר לטעות בה: השרתים קורסים, המכירות רצות מהר מהגיוס,
והצוות בקושי מספיק לתמוך בלקוחות~\cite{andreessen2007pmf}.

\section{מאמצים מוקדמים ולקוחות ראשונים}

לפני \en{PMF}, היעד אינו צמיחה אלא \emph{למידה}. לקוחות מוקדמים
\term{(early adopters)} מוכנים לסבול מוצר לא מושלם משום שהכאב שלהם חריף; הם
מקור המשוב היקר ביותר. בשלב זה עדיף ``לעשות דברים שאינם ניתנים להרחבה'':
לגייס לקוחות ידנית, לתמוך בהם אישית, וללמוד מכל שיחה.

\section{מדידת ההתאמה}

כיצד יודעים שהגענו? סימנים כמותיים כוללים שיעור שימור גבוה
\term{(retention)} שמתייצב ואינו דועך לאפס, ושיעור גבוה של משתמשים המצהירים
שהיו ``מאוכזבים מאוד'' אם המוצר ייעלם. כל עוד אין \en{PMF}, השקעה בשיווק נרחב
היא בזבוז: היא שופכת מים לדלי מחורר. רק לאחר שההתאמה מושגת, עוברים לשלב הצמיחה
ולכלכלת היחידה הנדונה בפרק~\ref{chap:economics}.
